\subsection{Good inert places and translated theta sections}
\label{subsec:cm-models}
\hypertarget{source-cm-models}{}

The formal group computation applies to the period kernel only if the translated theta sections extend on the corresponding integral CM model.

Let \(K=\mathbb Q(\sqrt{-D})\), where \(D>0\) is squarefree
and \(D\equiv3\bmod4\). The period lattice \(\Gamma\subset\C\)
is stable under \((1+\sqrt{-D})/2\). Thus its multiplier order
is the full ring of integers of \(K\). Fix \(p\ge5\) inert in
\(K\) and an abstract field isomorphism \(\C\simeq\C_p\).
All analytic identities below are transported on Taylor
coefficients. Put \(R=W(\overline{\mathbb F}_p)\), the integers
of the completed maximal unramified field, and let \(\sigma\)
be its Witt Frobenius.

\begin{proposition}[Good model and integral CM coordinates]\label{prop:good-inert-cm-model}
At the specified inert \(p\ge5\), a homothety of the complex
period lattice gives a good short model with coefficients in \(R\)
under the chosen transport. The
elliptic formal group has \(p\)-height two and admits an integral comparison with unit tangent
\[
 t(X)=l^{-1}(cL(X))\in R[\![X]\!],\qquad t'(0)=c\in R^\times.
\]

Nonzero prime-to-\(p\) torsion sections have integral affine
coordinates on the good model.
\end{proposition}

\begin{proof}
\proofstep{Obtain good reduction with full level three.}

Equip the elliptic curve with a
full level-three basis. The main theorem of complex multiplication
gives the moduli point with this basis over the ray class field
of modulus \(3\) over \(K\): under reciprocity the local units
preserve the ideal lattice and the units congruent to \(1\) at
\(3\) preserve the basis modulo \(3\). The full basis removes
automorphisms, so the curve with level structure descends there.
This field is unramified at \(p\). The singular modulus is
integral, so the curve has potentially good reduction. Its
inertia image on the \(3\)-adic Tate module is finite and
congruent to the identity modulo \(3\); it is trivial, since
the matrix logarithm on \(1+3M_2(\mathbb Z_3)\) has the
matrix exponential as inverse, and a torsion matrix there
has logarithm zero.
Indeed a torsion element \(A\in1+3M_2(\mathbb Z_3)\) of order \(n\) satisfies
\[
  0=\log(A^n)=n\log A\quad\Longrightarrow\quad \log A=0\quad\Longrightarrow\quad A=I .
\]

The good reduction criterion applies
since \(3\ne p\). A good Weierstrass model can be shortened
over \(R\) because \(2,3\) are units. Matching an invariant
differential scales \(g_2,g_3\) by the fourth and sixth
powers of a single factor.
Thus a change of scale by \(u\) preserves the quotient defining \(j\), since
\[
  (u^4g_2)^3=u^{12}g_2^3,\qquad (u^6g_3)^2=u^{12}g_3^2.
\]
 CM endomorphisms extend over the
valuation ring by the extension property of abelian schemes;
they are defined over the moduli field containing \(K\),
as can also be seen from their faithful action on the tangent
line. \cite{cm-reciprocity,serre-tate}

\proofstep{Determine the height and lift the formal isomorphism.}

The formal completion is a one-dimensional formal
\({\cal O}_{K_p}\)-module. The action of \({\cal O}_K\) extends
by \(p\)-adic continuity: in logarithmic coordinates the
coefficients of \([a](t)=l^{-1}(a l(t))\) at each degree are
polynomials in \(a\), and \(R\) is closed. The special fibre
of the elliptic formal group has \(p\)-height \(h=1\) or \(2\).
The first nonzero term of \([p]\) in the residue field has
degree \(p^h\). Let its leading coefficient be \(A\ne0\).

In the residue
field the leading terms of \([p]\circ[a]\) and
\([a]\circ[p]\) agree, so
\[
                    A\bar a^{p^h}= \bar a A
                            \quad(\bar a\in\mathbb F_{p^2}).
\]
Thus \(2\mid h\), and \(h\in\{1,2\}\) gives \(h=2\). The height as
a formal \({\cal O}_{K_p}\)-module is one. Over an algebraically
closed residue field such modules of the same height and
structural embedding are isomorphic, and the deformation ring
of a height-one formal module is just \(R\). The model with
logarithm \(L\) in \eqref{eq:integral-log} has this height:
its \([p]\)-series has linear coefficient \(p\), no terms of
degrees \(2,\ldots,p^2-1\), and coefficient \(1-p^{p^2-1}\)
at degree \(p^2\), by comparison in \(L([p]X)=pL(X)\).

In particular the first surviving term on reduction is
\[
  [p](X)\equiv X^{p^2}\pmod{(p,X^{p^2+1})}.
\]
It has the integral action of the specified quadratic integers
by \eqref{eq:integral-log}. Lift an isomorphism of the two
reductions by the height-one deformation theorem. Comparing
logarithms gives an integral isomorphism \(t(X)=l^{-1}(cL(X))\)
with \(c\in R^\times\) its tangent map. Nonzero prime-to-\(p\)
torsion on the good curve has integral affine coordinates in
\(R\): the corresponding torsion group scheme is finite
\'etale over the strictly henselian ring, so its distinct
sections reduce distinctly, including the zero section.
\cite{drinfeld-formal,serre-tate}

\end{proof}

Normalize the analytic invariants by the good-model homothety.
For the quasi-period homomorphism of the Weierstrass zeta
function write \(\eta(\omega)=s_2\omega+B\bar\omega\) on
\(\Gamma\).

\begin{proposition}[The integral norm relation]\label{prop:good-cm-norm}
For the normalized good model the quadratic correction \(s_2\)
belongs to \(R\). For every quadratic unit \(a\) at \(p\), its action in the good formal coordinate satisfies the congruence \eqref{eq:cm-norm-relation}
\[
 q([a]X)\sim a\sigma(a)q(X).
\]

\end{proposition}

\begin{proof}
\proofstep{Compare the isogeny divisor and its leading coefficient.}

The Legendre relation gives
\(B=\pi/\operatorname{area}(\C/\Gamma)\).
Then \(\varsigma_*(z)=\varsigma(z)e^{-s_2z^2/2}\) obeys
\[
 \varsigma_*(z+\omega)=\chi(\omega)
         \exp\big(B(z\bar\omega+|\omega|^2/2)\big)
         \varsigma_*(z),\qquad \chi(\omega)\in\{1,-1\}.
\]
For \(0\ne a\in{\cal O}_K\) let \(n=a\bar a\). On the complex
elliptic curve the following identity holds:
\begin{equation}\label{eq:cm-isogeny}
 \frac{\varsigma_*(az)^2}{\varsigma_*(z)^{2n}}
       =a^2\prod_{P\in\ker[a]\setminus\{0\}}
                                      (\wp(z)-\wp(P)).
\end{equation}
The left side is elliptic since the Hermitian exponent scales
by \(a\bar a\) and the signs square to one. The two sides have
the same divisor.

For a point with \(P=-P\) the factor on the
right has a double zero there; otherwise the two factors for
\(P,-P\) provide the double zero. Comparison of the first
Laurent coefficients at zero gives \(a^2\). Comparison at
the next even power gives
\[
                 (a^2-n)s_2=\sum_{0\ne P\in\ker[a]}\wp(P).
\]
Take \(a=\sqrt{-D}\). All these points have order dividing
\(D\), and \(a^2-n=-2D\) is a unit at \(p\). The good torsion
calculation therefore puts \(s_2\) in \(R\).

\proofstep{Use prime-to-\(p\) torsion and approximate the local units.}

For \(a\) a unit at \(p\), express \eqref{eq:cm-isogeny} in
the integral parameter \(t\). With \(A(t)=t^2\wp(l(t))\)
the result is
\[
 \exp\{-2(q_E([a]t)-nq_E(t))\}
  =\left(\frac{a t}{[a]t}\right)^2
                   \prod_{0\ne P\in\ker[a]}(A(t)-\wp(P)t^2).
\]
The right side is a principal integral unit. The factor
\([a]t/(a t)\) is integral with constant one by the extension
of the CM action, and all kernel points have prime-to-\(p\)
order.

Taking a square root is integral since \(p\) is odd.
Conjugate by \(t(X)\) and absorb the unit \(t(X)/(cX)\) as
in the preceding subsection. The norm relation
\eqref{eq:cm-norm-relation} follows for \(a\in{\cal O}_K\)
of prime-to-\(p\) norm. Such elements are dense in
\({\cal O}_{K_p}^\times\): approximate the two coordinates
in the integral basis \(1,(1+\sqrt{-D})/2\) modulo \(p^j\).
The first approximation modulo \(p\) already ensures the
norm is prime to \(p\).
Indeed, for the approximants \(a_j\) to a unit \(u\),
\[
  a_j\equiv u\pmod {p^j}\quad\Longrightarrow\quad a_j\bar a_j\equiv u\bar u\not\equiv0\pmod p .
\]
 On this quadratic field \(\sigma\)
is conjugation. Each coefficient of the exponential of the
norm difference is continuous in \(a,\sigma(a)\); its value
on all units is integral by closedness. This supplies the
comparison with all scalars in the trace calculation.

\end{proof}

Use the Heisenberg translation
\[
 (U_P f)(z)=f(z+P)\exp\{-B(z\bar P+|P|^2/2)\}.
\]
\begin{proposition}[Integral gauged torsion translates]\label{prop:good-cm-torsion-translates}
For each prime-to-\(p\) translate \(P\) of the normalized even
section \(\vartheta_0\), write \(G_P(t)=e^{q_E(t)}(U_P\vartheta_0)(l(t))\). The integrality and primitivity conclusions are
\[
 G_P(t)\in R[\![t]\!],\qquad G_P(t)\bmod\mathfrak m_R\ne0.
\]
 The same assertion holds in the
integral coordinate \(X\) with its gauge \(q\).
\end{proposition}

\begin{proof}
\proofstep{Obtain the addition recurrence.}

In \(t\) coordinates write \(G_P(t)=e^{q_E(t)}
 (U_P\varsigma_*)(l(t))\). In particular \(G_0=t\).
For \(P\) a period this is \(t\) times a sign. If \(Q\)
represents a nonzero torsion point, \(C=G_Q(0)\ne0\).
Sigma addition, with all quadratic exponential factors
cancelled, gives for \(P\) nonzero modulo periods
\begin{equation}\label{eq:torsion-difference}
      G_{P+Q}G_{P-Q}
            = C^2\{\wp(Q)-\wp(l(t)+P)\}G_P^2 .
\end{equation}
\proofstep{Determine the exact special-fibre orders.}

Take \(N\) the exact order of \(Q\), prime to \(p\). For
\(1\le j<N\) write \(f_j=\wp(Q)-\wp(l(t)+jQ)\).
It is an integral unit times \(t^{b_j}\), with
\(b_j={\bf1}_{j=1}+{\bf1}_{j=N-1}\). To check this, translation
by the integral section \(jQ\) identifies the completed local
ring there with \(R[[t]]\); \(x\) is regular on its neighbourhood.
The divisor \(x-x(Q)\) on the smooth Weierstrass curve has
zeros exactly at the sections \(Q,-Q\), counted with
multiplicity. On the special fibre its zeros have the same
description. Reduction injects this prime-to-\(p\) torsion
group; at order two the double zero is exact since the cubic
has distinct roots and \(p>2\). Thus the first nonzero
coefficient at the indicated order is a unit.

\proofstep{Recover the scalar and the integral series.}

Repeatedly using \eqref{eq:torsion-difference} gives
\[
 G_{NQ} G_0^{\,N-1}
       =G_Q^{\,N}\prod_{j=1}^{N-1}(C^2 f_j)^{N-j}.
\]
This follows by taking successive quotients \(G_{(j+1)Q}/G_{jQ}\)
in the Laurent series field. The total vanishing order of the
product of \(f_j\)'s is \(N\), including \(N=2\). The left side
is a unit scalar times \(t^N\). Cancel powers of \(t\) and
compare constants: 
\[
  |C|_p^{N^2}=1\quad\Longrightarrow\quad v_p(C)=0.
\]
 Hence \(G_Q^N\) is
an integral unit, and \(G_Q/C\) itself is one by the binomial
series with exponent \(1/N\). For \(Q\) representing zero the
formula \(G_Q=\pm t\) gives integrality directly.

Fix a
nonzero half-period \(Q_0\). Normalize the even translate
\(\vartheta_0=U_{Q_0}\varsigma_*/(U_{Q_0}\varsigma_*)(0)\).
Composing translations by \(P,Q_0\) differs from translation
by \(P+Q_0\) by a root of unity for torsion arguments, since
the alternating lattice pairing takes rational values on
division points. The denominator at \(Q_0\) is a unit by the
order-two case. Thus for any prime-to-\(p\) translate \(P\),
\(e^{q_E(t)}(U_P\vartheta_0)(l(t))\) is integral with some
unit coefficient. Substitution of the integral isomorphism
\(t(X)\) and the change \(q_E\mapsto q\) preserve this
assertion.

\end{proof}

\begin{auxlemma}[Extracting a coefficient or a derivative from disjoint copies]
\label{lem:cm-disjoint-tensor}
Let \(F\) be a valued field of characteristic zero and
\(H(X)=\sum_\alpha a_\alpha X^\alpha\in F[[X_1,\ldots,X_n]]\).
For any \(m\ge1\), take disjoint sets of variables \(X^{(1)},\ldots,X^{(m)}\)
and set \(T=\prod_{j=1}^m H(X^{(j)})\). For a multi-index \(\nu\),
\[
 \left[\prod_j(X^{(j)})^\nu\right]T=a_\nu^m,\qquad
 \left.\left(\prod_j\partial_{X^{(j)}}^\nu\right)T\right|_0
         =\bigl((\partial_X^\nu H)(0)\bigr)^m .
\]
Integrality of the indicated coefficient or derivative of \(T\)
therefore gives integrality of the corresponding coefficient or derivative of \(H\).
\end{auxlemma}

\begin{proof}
In the indicated monomial each set of variables has its own exponent \(\nu\),
so precisely one product of coefficients occurs. Multiplying each factor
by \(\nu!=\prod_i\nu_i!\) gives the derivative identity. If a nonzero
element has integral \(m\)-th power, then \(m\) times its valuation is
nonnegative. The zero reading is integral as well.
\end{proof}

\begin{proposition}[Integral correction in an inert CM chart]\label{prop:inert-cm-chart}
For the full-order curve at an inert prime \(p\ge5\),
use the good formal model after homothety. Suppose that, after an integral unitary change of orthogonal frame,
the fourth tensor power of the translated theta section is a finite
integral linear combination of products of the one-variable
translates of Proposition~\ref{prop:good-cm-torsion-translates}.
For formal coordinates \(\mathbf X\) on that chart, the corrected series satisfies
\[
 e^{\sum q}\Theta\in R[\![\mathbf X]\!].
\]
This includes \(\Theta(0)=0\). Here \(q\) is the quadratic correction
series from Section~\ref{subsec:curve-exponents}.
\end{proposition}

\begin{proof}
\proofstep{Transport the gauge by the integral unitary frame.}

To see where these one-variable units enter the orthogonal
chart, let \(\gamma\) be an integral unitary change of frame
over the quadratic integers at \(p\). Its inverse conjugate
transpose is integral. On \(4n\) coordinates it acts by
\(Z\mapsto L^{-1}(\gamma L(Z))\) componentwise. By the CM
norm relation and additivity of \(b\), the sum of \(q\) over
all coordinates is invariant modulo \(\sim\): the diagonal
contractions of \(\gamma^*\gamma\) are one and the off-diagonal
ones vanish.

For the contracted coefficients the two cases read
\[
  \sum_i\overline{\gamma_{ij}}\gamma_{ik}=(\gamma^*\gamma)_{jk}=\delta_{jk}.
\]
 Suppose in this frame the fourth tensor of the
translated theta section has a finite expansion in products
of the one-variable translates above, with integral coefficients.
After multiplication by the exponential of the summed \(q\)
each product is integral. Undo the unitary substitution and
absorb the principal unit from the invariant sum. The result
is that the fourth tensor of \(e^{\sum q}\Theta\) is integral.
\proofstep{Recover each coefficient from its fourth power in disjoint variables.}

Lemma~\ref{lem:cm-disjoint-tensor}, with \(m=4\) and 
\[
  H=e^{\sum q}\Theta,
\]
now gives integrality coefficient by coefficient, including at a zero
of the section.

The global scalar and its local standard frames are constructed in Section~\ref{subsec:theta-normalization}.
\label{cor:cm-models-1}
\end{proof}
